% Generated draft: requires mathematical and source review.
\documentclass[11pt]{article}
\usepackage[margin=25mm]{geometry}
\usepackage{fontspec}
\setmainfont{Latin Modern Roman}
\usepackage{amsmath,amssymb,mathtools,unicode-math}
\setmathfont{Latin Modern Math}
\usepackage{enumitem,needspace,xurl,hyperref}
\setlength{\parindent}{0pt}
\setlength{\parskip}{5pt}
\setlength{\emergencystretch}{3em}
\newcommand{\N}{\mathbb{N}}
\newcommand{\Z}{\mathbb{Z}}
\newcommand{\Q}{\mathbb{Q}}
\newcommand{\R}{\mathbb{R}}
\newcommand{\C}{\mathbb{C}}
\newcommand{\F}{\mathbb{F}}
\newcommand{\A}{\mathbb{A}}
\newcommand{\PP}{\mathbb{P}}
\newcommand{\E}{\mathbb{E}}
\newcommand{\eps}{\varepsilon}
\newcommand{\dd}{\,\mathrm{d}}
\begin{document}
\section*{UnsolvedMath ID: 30003861}
\textbf{Source problem code:} catalogue entry
\subsection{Definitions and mathematical statement}
Write \(D_Y\) for the boundary divisor and set
\[
U=Y\setminus\operatorname{Supp}(D_Y).
\]
A simple normal-crossings boundary is a reduced divisor whose components are smooth and meet locally as distinct coordinate hyperplanes. The usual linear-triviality condition for a log Calabi--Yau pair is \(K_Y+D_Y\sim 0\); the supplied record does not specify its convention for this condition. For an \(n\)-dimensional pair, a common meaning of maximal boundary is the existence of a zero-dimensional boundary stratum, where \(n\) components meet. This convention is not explicitly supplied either. The hypothesis that \(D_Y\) supports an ample divisor asserts the existence of an ample divisor supported on the boundary; the allowed divisor coefficients and whether its support must equal the entire boundary are unspecified.

\emph{Intended question, with unresolved qualifications.} For every simple normal-crossings Calabi--Yau pair \((Y,D_Y)\) with maximal boundary and with \(D_Y\) supporting an ample divisor, does the Gross--Hacking--Keel curve-counting algebra construction produce a flat family of affine log Calabi--Yau varieties such that its torus quotient depends only on \(U\) and is a mirror family of \(U\)? Here dependence only on \(U\) expresses the proposed independence from the choice of compactifying pair \((Y,D_Y)\).

This is a statement of the supplied intended target, not a fully specified conjecture. The record supplies neither the ground field and dimension restrictions nor the curve-counting algebra, its coefficient ring, the family base, the acting torus, the quotient construction, or a mathematical criterion for being a mirror family. It also does not specify the equivalence under which the quotient is to depend only on \(U\). Consequently these data cannot be replaced by a precise quantified construction from the supplied evidence alone.
\subsection{Short English statement}
The question asks whether the Gross--Hacking--Keel curve-counting construction gives a flat affine log Calabi--Yau family whose torus quotient is an intrinsic mirror of the complement \(U\). Essential construction data and conventions remain unspecified. The supplied review also reports an incomplete source extraction, despite the record's separate claim that the statement was verified.
\subsection{Sources}
\begin{itemize}
\item Supplied local UnsolvedMath registry record, permanent ID \(30003861\), title ``Mirror Families from Log Calabi--Yau Pairs,'' exact source problem code \texttt{catalogue entry}, in the Oberwolfach Reports Open Problems collection. The proposed year is recorded as 2018; no proposer is named. The construction is attributed to Gross--Hacking--Keel, without a separate construction reference.
\item Supplied report citation: \emph{Subgroups of Cremona Groups} (2019), Oberwolfach Reports, pp.~1685--1743, DOI \texttt{10.4171/owr/2018/28}. \url{https://doi.org/10.4171/owr/2018/28}. The citation's 2019 date and the report's 2018 identifier are retained as supplied.
\item Supplied literature source: ``Oberwolfach Report 28/2018.'' \url{https://ems.press/content/serial-article-files/46750}.
\item Literature assessment dated exactly \texttt{2026-08-21}: the supplied curation says the extracted record ends mid-conjecture and lacks a complete mathematical claim; it assigns the label ``open'' under its queue rule because a solved status could not be verified. This is supplied dated curation, not an independently verified literature assessment or proof of openness. The report and cited claims have not been independently consulted for this draft.
\item The record separately labels the statement \texttt{corrected\_verified} and reports a manual semantic, typographic, and LaTeX review asserting self-containment and preservation of the target. That supplied assessment does not resolve the conflicting incomplete-extraction qualification.
\end{itemize}
\end{document}
